\scanpage{41}
% archive foliation: 40
\uline{Dém} On a déjà vu l'équivalence de $(1)$ à $(3)$, et \deleted{\unsure{(3)}} \inserted{\unsure{ce|que}} implique $4$ (i.e. $e_v(\unsure{0|1})f=fe_u(\unsure{0|1})$). Prouvons $(4)\longrightarrow(3)$. Or $(4)$ s'écrit
\[
\begin{gathered}
\underbrace{(\exp v)\circ f-f\circ\exp u}_{\text{''}}=0\\
\sum_{n\geq1}\frac1{n!}\underbrace{(v^nf-fu^n)}_{\delta_n\ \mathrm{déf}}
\end{gathered}
\]
Or on a (réc. sur $n$)
\[
\delta_n=\unsure{v|u}^{n-1}\delta+v^{n-2}\delta u+\cdots+v\delta u^{n-2}+\delta u^{n-1}
\]
donc $(4)$ signifie
\[
\delta+\frac12(v\delta+\delta u)+\frac1{3!}(v^2\delta+v\delta u+\delta u^2)+\cdots=0
\]
\margin{(somme \uline{finie})}
On peut l'écrire
\[
\delta=\sum_{i+j\geq1}c_{ij}v^i\delta u^j
\qquad\begin{array}{c}(c_{ij}\in K)\\\text{''}\\\dfrac1{(i+j)!}\end{array}
\]
On met $U=\operatorname{End}_K(M)$
\[
\Theta,\Theta':U\longrightarrow U\qquad K\text{-endom}
\]
\[
\Theta(x)=vx,\qquad\Theta'(x)=xu
\]
Donc $\Theta,\Theta'$ sont nilpotents ($v,u$ l'étant) et commutent, et on aura
\[
\delta=\sum_{i+j\geq1}c_{ij}\Theta^i\Theta'{}^j\delta=\rho\delta
\]
\[
\rho=\sum_{i+j\geq1}c_{ij}\Theta^i\Theta'{}^j\in\operatorname{End}_K(U)
\]
Or $\rho$ est \uline{nilpotent} ($\Theta,\Theta'$ l'étant et commutant).

Or $\delta=\rho\delta$ implique par récurrence
